\section{Classification Difficulty and Population Information}
\label{sec:tail-identifiability}

\subsection{Efficient Information and Local Classification Difficulty}
\label{sec:classification-difficulty}

Parameter ambiguity need not imply classification difficulty: probabilities $0.90$ and $0.95$ prescribe the same label. A lower bound requires statistically similar alternatives whose optimal decisions disagree on enough probability mass. Define classification risk $R_P(g)=P\{g(X)\ne Y\}$, Bayes risk $R_P^\star=\inf_g R_P(g)$, and excess risk $\mathcal E_P(g)=R_P(g)-R_P^\star$. For a specified local family, let $\mathfrak R_e(n)$ denote minimax expected excess classification risk based on $n$ independent observations from environment $e$. Write $\KL(P\Vert Q)$ for the Kullback--Leibler divergence.

\begin{theorem}[Local classification difficulty]
\label{thm:local-classification-difficulty}
Fix $\lambda_e>0$ and a unit weakest efficient-information direction $u_e$. For every $0<t\le t_{0,e}$, suppose the local family contains $P_{e,+t},P_{e,-t}$ with task coordinates $\pm tu_e$ and nuisance coordinates allowed to vary. Assume, for constants $K_e,c_e,t_{0,e}>0$ and $q>0$,
\begin{equation}
\KL(P_{e,+t}\Vert P_{e,-t})\le K_et^2\lambda_e,
\qquad
\inf_g\{\mathcal E_{P_{e,+t}}(g)+\mathcal E_{P_{e,-t}}(g)\}\ge c_et^q.
\label{eq:local-classification-assumptions}
\end{equation}
Then, for a universal constant $c_0>0$,
\begin{equation}
\mathfrak R_e(n)\ge c_0c_e
\min\{t_{0,e}^q,(nK_e\lambda_e)^{-q/2}\}.
\label{eq:local-classification-lower-bound}
\end{equation}
\end{theorem}

The KL bound limits distinguishability; decision separation supplies the classification cost. Both must hold over radius $t_{0,e}$. With quadratic separation, the interior rate is $c_e/(nK_e\lambda_e)$, capped by the local decision scale. Appendix~\ref{app:theory-classification} proves the testing reduction; the next example calculates both conditions and a matching upper bound.

\begin{samepage}
\subsection{A Worked Classification Mechanism}
\label{sec:classification-mechanism}
Let $X\sim\operatorname{Unif}[-1,1]$, $0<\delta\le1/2$, $|\theta|,|\mu|\le\delta/8$, and
\begin{equation}
\eta_{\theta,\mu}(x):=P_{\theta,\mu}(Y=1\mid X=x)
=\tfrac12+\tfrac{x}{4}+\theta+\mu\mathbf1\{|x|>\delta\}.
\label{eq:boundary-region-family}
\end{equation}
Here $\delta$ is boundary-region half-width and probability mass. The nuisance reproduces a task shift outside that region; the Bayes rule is $g_\theta^\star(x)=\mathbf1\{x\ge-4\theta\}$ throughout the family. At $(\theta,\mu)=(0,0)$,
\begin{equation}
J_{TT}=4\log3,
\qquad F_\delta=4\log\frac{2+\delta}{2-\delta}\asymp\delta.
\label{eq:boundary-region-information}
\end{equation}
The alternatives $(t,-t)$ and $(-t,t)$ agree outside $|X|\le\delta$, have decision separation $4t^2$, and satisfy $\KL(P_{+t}\Vert P_{-t})\le5F_\delta t^2$ for $0<t\le\delta/8$. Theorem~\ref{thm:local-classification-difficulty} and a matching estimator give minimax risk over this parameter rectangle, uniformly in $n\ge1$ and $\delta$,
\begin{equation}
\mathfrak R_\delta(n)\asymp\min\{\delta^2,(n\delta)^{-1}\}.
\label{eq:boundary-region-minimax-rate}
\end{equation}
Approximately $n\delta$ observations distinguish the task effect from nuisance variation where the decision boundary can move. The $\delta^2$ cap retains the shrinking local decision scale. This statistical nuisance preserves Bayes labels at fixed $\theta$; it does not model the physical background-swap operation. Appendix~\ref{app:classification-mechanism} gives the calculation and a finite-sample upper bound.\par
\end{samepage}

\subsection{Information Tails Across Populations}
\label{sec:collapse-geometry}

Population aggregation describes how often weak task information occurs; classification still retains the radius and separation factors above. The inherited geometric law is:
\begin{theorem}[Geometric tail law {\citep{mansouri_arai_2026_fisher_glasses}}]
\label{thm:inverse-information-geometry}
Under the uniform regularity and localization assumptions in Appendix~\ref{app:theory-tails}, let the Fisher-null set have codimension $c$, information vanish to order $2m$, and nearby density have order $\nu>-c$. Then, for some $C>0$,
\begin{equation}
\Pr(\lambda(z)^{-1}>x)\sim Cx^{-\beta},
\qquad \beta=\frac{c+\nu}{2m}.
\label{eq:inverse-information-tail-exponent}
\end{equation}
\end{theorem}
Quadratic collapse with positive density gives $\beta=c/2$; the inverse mean is finite exactly when $\beta>1$. Appendix~\ref{app:theory-tails} contains the calculation and Figure~\ref{fig:geometry-tails}.

\begin{samepage}
\subsection{The Fisher-Glass Certificate}
\label{sec:fg-certificate}

Upper-tail conditional value-at-risk, $\CVaR_\alpha$, averages the largest $1-\alpha$ fraction of a quantity, with fractional weighting at atoms \citep{rockafellar2000optimization,rockafellar2002conditional}. We apply the Fisher-glass certificate \citep{mansouri_arai_2026_fisher_glasses} to classification environments $e\sim\Pi$, using the weakest task information $\lambda_e$ and stabilization $\varepsilon$:
\begin{equation}
L_e^{(\varepsilon)}=(\lambda_e+\varepsilon)^{-1},
\qquad
I_\alpha^{\FG}=\bigl[\CVaR_{\alpha,e\sim\Pi}(L_e^{(\varepsilon)})\bigr]^{-1},
\quad 0<\alpha<1.
\label{eq:fisher-glass-certificate}
\end{equation}
\end{samepage}
A larger certificate means a smaller inverse-information tail, expressed in information units. Stabilization $\varepsilon>0$ caps inverse information at $1/\varepsilon$; it does not remove the underlying null geometry. Under Theorem~\ref{thm:inverse-information-geometry}'s power-law regime, the limit $\varepsilon\downarrow0$ is positive for $\beta>1$ and vanishes for $\beta\le1$.

\begin{proposition}[Information coverage]
\label{prop:information-coverage}
For $\rho>\varepsilon$,
\begin{equation}
I_\alpha^{\FG}\ge\rho
\quad\Longrightarrow\quad
\Pr_{e\sim\Pi}(\lambda_e\ge\rho-\varepsilon)\ge\alpha.
\label{eq:information-coverage}
\end{equation}
\end{proposition}

The certificate establishes an absolute information floor on an $\alpha$-fraction of the declared environment distribution; classification consequences retain the local radius and decision separation. Appendix~\ref{app:theory-certificate} gives the proof. With four equally weighted groups and $\alpha=0.95$, it reduces to the minimum stabilized information, so that real-data audit does not isolate a benefit from tail aggregation.
